\section{Problem Setting}\label{sec:setting}
\label{sec:prelim}
\label{sec:problem}

In the cyclic partially observable Markov decision process (POMDP), the latent
position is $s_t\in\Z_k$, with known cycle size $k$ and positions numbered from
$0$ to $k-1$ \citep{kaelbling1998pomdp}. The increment $j_t$ gives
$s_t=(s_{t-1}+j_t)\bmod k$. An observation $o_t\in\R^n$ can contain a cue about
$j_t$ and a pointer reading $y_t$ of the absolute position $s_t$. The agent can
choose an action $a_t$ after observing $o_t$. In control, actions can affect the
next increment or request a new pointer reading. The belief is the posterior
probability of each position,
\[
b_t(s)=\Pr(s_t=s\mid o_{1:t}),
\]
with past actions included in the observation history for controlled sequences.
A learned policy updates a memory $h_t=f_\phi(o_t,h_{t-1})$ and assigns action
probabilities $\pi_\phi(a\mid h_t)$.

We evaluate belief estimates beyond the training lengths using the accuracy of
the most probable position and the negative log likelihood (NLL) of the true
position, $-\log b_t(s_t)$. The fixed shared planner selects actions from each
estimator's belief to maximize expected net utility, defined as the success
indicator minus the costs of turns and readings
(Section~\ref{sec:regimes}).

\myparagraph{Counting tasks.}
In the \textbf{mod-$k$} task, each of the $L$ observations encodes a binary event
bit $e_t\in\{0,1\}$. A single decision at the final step earns $+1$ for the
correct count $a^\star=\big(\sum_{t=1}^{L} e_t\big)\bmod k$ and $-1$ otherwise.
Parity is the case $k=2$. Models train on a bounded range of lengths and are
evaluated on longer ones. On clean inputs, the programmed counter
$c_t=(c_{t-1}+e_t)\bmod k$ is an exact oracle. The experiment tests whether a
policy can learn the count from terminal reward. Two control tasks remove the
cyclic structure. In the \textbf{cue} task a single informative bit appears at
$t=1$. In the \textbf{monotone} task the agent reports whether the count exceeds
$L/2$. Sections~\ref{sec:transfer} and~\ref{sec:feedback} describe the valve
control and missing-event tasks.
